\documentclass[11pt]{article}

% ---- xelatex + Korean ----
\usepackage{kotex}
\usepackage[margin=1in]{geometry}
\usepackage{amsmath,amssymb,amsthm}
\usepackage{mathtools}
\usepackage{booktabs}
\usepackage{array}
\usepackage{xcolor}
\usepackage{enumitem}
\usepackage[colorlinks=true,linkcolor=blue!55!black,urlcolor=teal]{hyperref}

\newcommand{\code}[1]{\texttt{\small #1}}
\newcommand{\pth}{p_\theta}
\newcommand{\pphi}{p_\phi}
\newcommand{\Qb}{\bar{Q}}
\newcommand{\xz}{x_0}
\newcommand{\xt}{x_t}
\newcommand{\xpre}{x^{<b}}
\newcommand{\bLM}{\mathcal{M}}
\DeclareMathOperator{\Cat}{Cat}
\DeclareMathOperator{\KL}{KL}
\DeclareMathOperator{\softmax}{softmax}
\DeclareMathOperator{\CE}{CE}
\DeclareMathOperator*{\E}{\mathbb{E}}

\theoremstyle{definition}
\newtheorem{remark}{Remark}

\title{\bfseries Block Diffusion 학습 Loss 정식화\\[2pt]
\large 확산 모델 학습 loss (mask / graph 커널) + Discriminator·Classifier 학습 loss}
\author{Block-Diffusion-OD-Planning \quad (v2 backbones) \\[-1pt]
\small 자매 문서: \code{BD\_BLOCK\_GUIDANCE\_MATH.pdf}, \code{BD\_DCBG\_MATH.pdf}}
\date{2026-07-21}

\begin{document}
\maketitle

\begin{abstract}
\noindent
본 문서는 학습에 쓰이는 모든 loss 를 수식으로 정리한다:
(1) block-diffusion 확산 모델 학습 loss —흡수(mask) 커널의 SUBS 가중 CE 와
균일(graph) 커널의 D3PM ELBO(+보조 CE·연결성 정칙화)—,
(2) guidance 용 두 판별기의 학습 loss —IW discriminator(clean 부분경로 BCE)와
D-CBG classifier(노이즈 상태 BCE). 각 수식에 구현 위치(\code{models\_seq/bd\_models.py},
\code{bd\_trainer.py}, \code{train\_bd\_disc.py}, \code{train\_dcbg\_classifier.py})를 병기한다.
\end{abstract}

% =====================================================================
\section{Canvas 와 loss mask (길이 처리)}
% =====================================================================
경로 $[\,\text{ori},v_1,\dots,v_L{=}\text{dst}\,]$ 를 canvas
$x_0=[\text{dst},\text{ori},v_1,\dots,v_L{=}\text{dst},\langle\text{END}\rangle,\dots\mid \text{PAD},\dots]$
로 만든다(\code{build\_canvas}). 학습 대상 집합 $\bLM$ (loss mask)은:
\begin{itemize}[nosep,leftmargin=1.4em]
  \item \textbf{학습됨}: 경로 본문 $v_1..v_L$ + dst 가 속한 블록의 $\langle\text{END}\rangle$ tail.
  \item \textbf{제외}: dst 블록 이후의 \textbf{PAD 블록 전체} (노이즈·loss 대상 아님).
  \item 조건 dropout 된 표본은 prefix $[\text{dst},\text{ori}]$ 도 노이즈+학습(무조건 생성용): $\bLM[\neg\text{cond},{:}\text{pfx}]{=}\text{True}$
        (\code{bd\_models.py:435, 498}).
\end{itemize}
블록별 시각 $t_b$ 를 뽑아 토큰별 $t_{\text{tok}}=t_b$ (블록 내 공유)로 확장한다.

% =====================================================================
\section{확산 학습 loss — mask (흡수/SUBS) 커널}
% =====================================================================
블록별 $t_b\sim\mathcal U[10^{-3},1]$. 각 학습 위치 $i\in\bLM$ 는 확률 $t_i$ 로 [MASK] 로 치환:
\begin{equation}
\text{move}_i=\mathbf 1[u_i<t_i]\wedge\mathbf 1[i\in\bLM],\qquad
\xt^i=\begin{cases}\text{[MASK]}&\text{move}_i\\ \xz^i&\text{else}\end{cases}
\quad(\code{bd\_models.py:437--438}).
\end{equation}
SUBS 파라미터화: [MASK] 상태의 확률을 0 으로 두고($\text{logits}[\text{MASK}]{=}-\infty$) softmax,
$\log \pth(\xz^i\mid \xt,\xpre)=\big[\log\softmax(\text{logits}_i)\big]_{\xz^i}$
(\code{:451--453}). Loss 는 마스킹된 위치의 \textbf{$1/t$-가중 음의 로그가능도}:
\begin{equation}
\boxed{\;\mathcal L_{\text{mask}}
=\frac{1}{|\bLM|}\sum_{i}\ \text{move}_i\cdot\Big(-\tfrac{1}{t_i}\,\log \pth(\xz^i\mid \xt,\xpre)\Big)\;}
\quad(\code{bd\_models.py:456--458}).
\end{equation}
$1/t$ 가중은 MDLM/SUBS 흡수 ELBO 의 연속시간 가중이며, 정규화 분모는 loss-mask 크기 $|\bLM|$ 이다.
(bd3lms \code{diffusion.py:819--891} 대응.)

\subsection*{BD3-LM NELBO (식 8) 와의 대응}
BD3-LM 의 블록 확산 NELBO 는
\begin{equation}
-\log\pth(\mathbf x)\ \le\ \mathcal L_{\mathrm{BD}}
=\sum_{b=1}^{B}\E_{t\sim[0,1]}\,\E_q\,
\frac{\alpha'_t}{1-\alpha_t}\,\log\pth(\mathbf x^b\mid \mathbf x^b_t,\xpre).
\tag{BD3-LM 8}
\end{equation}
위 $\mathcal L_{\text{mask}}$ 가 이 식을 그대로 구현함을 세 부분으로 확인한다.

\textbf{(i) Weight 일치.} 코드는 선형(log-linear) 스케줄을 쓴다: 시각 $t$ 에서 마스킹 확률이
정확히 $t$ ($\text{move}_i=\mathbf 1[u_i<t_i]$, \code{bd\_models.py:437}) 이므로 미마스킹
확률은 $\alpha_t=1-t$. 따라서
\begin{equation}
\alpha'_t=-1,\qquad 1-\alpha_t=t
\quad\Longrightarrow\quad
\frac{\alpha'_t}{1-\alpha_t}=-\frac{1}{t},
\end{equation}
즉 식 8 의 weight 가 우리의 $-1/t$ 와 \emph{정확히} 일치한다(임의 스케줄 weight 의
선형 스케줄 특수화; cosine 등 다른 스케줄이면 형태가 달라진다).

\textbf{(ii) Token 분해는 정확(근사 아님).} SUBS 는 위치별 독립(mean-field) 복원이며
absorbing ELBO 에서는 \emph{마스킹된 위치만} 기여하므로 블록 재구성 항이 토큰합으로
분해된다:
\begin{equation}
\log\pth(\mathbf x^b\mid \mathbf x^b_t,\xpre)
=\sum_{i\in b}\mathbf 1[\text{move}_i]\,\log\pth(\xz^i\mid \xt,\xpre).
\end{equation}
그러므로 블록 항은 정확히 토큰별 $1/t$-가중 CE 의 합이고(\code{(nll*move).sum()},
\code{:456--458}), 식 8 의 $\sum_{b}$ 는 block-diffusion attention
(\code{get\_block\_diff\_mask}) 하의 전체 학습위치 합으로 실현된다($\xpre$ 는 clean 문맥).

\textbf{(iii) 정규화의 불편성.} \code{move}(마스킹 위치)로만 합하고 분모를
$|\bLM|$(전체 학습가능 위치)로 두는 것은, $\E_q[\text{move}_i]=t$ 에 의해
\begin{equation}
\E_q\!\Big[\textstyle\sum_{i}\text{move}_i\cdot\tfrac{1}{t}\big(-\log\pth^{\,i}\big)\Big]
=\sum_{i\in\bLM}\big(-\log\pth^{\,i}\big)
\end{equation}
이므로 $1/t$ 가중 $\times\,\tfrac{1}{|\bLM|}$ 가 토큰당 ELBO 의 \emph{불편(unbiased)}
추정이 된다 — 분모가 $|\text{masked}|$ 가 아니라 $|\bLM|$ 인 이유다.

\begin{remark}
이 대응은 \textbf{mask(흡수) 커널} 전용이다. graph(균일 CTMC) 커널은 식 8 이 아니라
아래 §의 이산 D3PM ELBO(KL${}+{}$보조 CE${}+{}$연결성 정칙화)로 별도 구현된다.
단, 이는 블록 \emph{내부} 커널의 차이일 뿐이며, 블록 \emph{사이}의 semi-AR
조건화($\xpre$)는 두 커널이 동일한 메커니즘으로 공유한다(다음 § 참고).
\end{remark}

% =====================================================================
\section{확산 학습 loss — graph (균일 CTMC/D3PM) 커널}
% =====================================================================
$\Qb_t$ =누적 전이행렬(\code{matrices}), $Q_t$ =한 스텝 전이(\code{Q}). 블록별 정수
$t_b\in\{1,\dots,T\}$ (uniform/early/cosine 샘플링, \code{bd\_models.py:479--491}).
확산 상태는 $q(\xt\mid\xz)=\Cat(\xz\Qb_t)$ 로 학습 위치에만 노이징(PAD clean):
\begin{equation}
\xt^i\sim \Cat\big([\Qb_{t_i}]_{\cdot,\xz^i}\big)\quad(\code{:504--506}),\qquad
\hat{x}_0=\softmax(\text{logits}_{:V_0})\quad(\code{:513}).
\end{equation}
\textbf{이전 블록 조건화($\xpre$)는 mask 커널과 동일한 메커니즘이다.} 모델 입력을
$x_{\text{in}}=[\,\xt\ \|\ \xz\,]$ 로 두고(noised 현재 $\|$ clean 전체) block-diffusion
attention 으로 넣는다(\code{get\_block\_diff\_mask}, \code{:508--511}). 이 attention 하에서
블록 $b$ 의 noised query 는 concat 된 clean $\xz$ half 의 \textbf{이전 블록들}에 attend 하므로,
$\hat{x}_0$ 및 예측 posterior 가 $\xpre$ 에 조건화된다 — 이것이 아래 $\pth(\cdot\mid\xt,\xpre)$
의 $\xpre$ 이다. 즉 graph 커널도 블록 \emph{사이}는 semi-AR(clean 이전 블록 조건),
블록 \emph{내부}만 균일 CTMC 확산이다.
\textbf{참(true) posterior} 와 \textbf{예측(pred) posterior} (D3PM):
\begin{align}
q(x_{t-1}{=}k\mid \xt,\xz)&\propto [Q_t]_{k,\xt}\,[\Qb_{t-1}]_{k,\xz}
&&(\code{:517--519}),\\
\pth(x_{t-1}{=}k\mid \xt,\xpre)&\propto [Q_t]_{k,\xt}\,\big[\Qb_{t-1}^{\top}\hat{x}_0\big]_{k}
&&(\code{:529--531}).
\end{align}
\begin{remark}[True posterior 의 차원별 독립성 — 이전 블록과 무관]
D3PM forward 는 차원별 독립이므로 ($q(\xt\mid\xz)=\prod_i q(\xt^i\mid\xz^i)$),
$\xz$ 를 알면 \textbf{true posterior 도 완전히 토큰(차원)별로 분해}되어 위치 $i$ 의
posterior 는 오직 $(\xt^i,\xz^i)$ 에만 의존한다
($q(x_{t-1}^i\mid\xt,\xz)\propto[Q_t]_{\cdot,\xt^i}\odot[\Qb_{t-1}]_{\cdot,\xz^i}$, \code{:517--519}).
즉 같은 블록의 다른 토큰과도, \emph{이전 블록과도 독립}이며 모델·attention 이 개입하지 않는다.
이전 블록 조건화 $\xpre$ 는 \textbf{오직 predicted posterior} 에만, 그것도 $\hat x_0$ 를 통해서
들어간다($\hat x_0$ 가 block-diffusion attention 으로 clean $\xpre$ 에 조건화, \code{:529--531}).
따라서 KL$(q\,\|\,\pth)$ 의 타깃 $q$ 는 $\xpre$-독립 고정분포이고, 블록 조건화는 확산
과정·true posterior 가 아니라 \textbf{denoiser($\hat x_0$ 예측기)에 대한 모델링 선택}일 뿐이다.
이런 의미에서 graph 커널 손실은 GDP(\code{seq\_models.py} Restorer)를 그대로 따르며, 블록
확산이 추가한 것은 (i) 블록별 독립 시간 $t_b$ 와 (ii) $\hat x_0$ 의 $\xpre$ 조건화, 이 둘뿐이다.
\end{remark}
per-token D3PM ELBO 항은 두 posterior 의 KL(masked mean):
\begin{equation}
\mathcal L_{\text{KL}}=\frac{1}{|\bLM|}\sum_{i\in\bLM}\KL\!\big(q(x_{t-1}\mid \xt,\xz^i)\,\|\,\pth(x_{t-1}\mid \xt,\xpre)\big)
\quad(\code{:537--538}).
\end{equation}
보조 $x_0$ 교차엔트로피와 연결성 정칙화(인접행렬 $A$ 위 bilinear, 연속 실정점쌍):
\begin{align}
\mathcal L_{\text{CE}}&=\frac{1}{|\bLM|}\sum_{i\in\bLM}-\log[\hat{x}_0^i]_{\xz^i}
\quad(\code{:539--540}),\\
\mathcal L_{\text{con}}&=-\frac{100}{Z}\sum_{\ell}\text{pair}_\ell\Big(
\hat{x}_0^{\ell}{}^{\!\top} A\,\log\hat{x}_0^{\ell+1}
+\hat{x}_0^{\ell+1}{}^{\!\top} A\,\log\hat{x}_0^{\ell}\Big)
\quad(\code{:545--552}).
\end{align}
\textbf{총 loss (ce-gating, \code{bd\_trainer.py:39--52})}: CE 가 임계 미만이면 KL 만, 아니면 가중합:
\begin{equation}
\boxed{\;\mathcal L_{\text{graph}}=
\begin{cases}
\mathcal L_{\text{KL}}, & \mathcal L_{\text{CE}}<60/16\ \text{또는}\ \lambda_{\text{ce}}{=}\lambda_{\text{con}}{=}0,\\[2pt]
\mathcal L_{\text{KL}}+\lambda_{\text{ce}}\mathcal L_{\text{CE}}+\lambda_{\text{con}}\mathcal L_{\text{con}}, & \text{그 외.}
\end{cases}\;}
\end{equation}
(\code{seq\_models.py} 의 Restorer.forward 수식을 블록-확산 attention 으로 이식.)

% =====================================================================
\section{판별기 학습 loss (1) — IW discriminator (clean 부분경로)}
% =====================================================================
IW disc 는 \emph{clean} 부분경로에서 ``실제 예외 데이터 vs negative''를 BCE 로 구분한다
(label smoothing $\epsilon{=}0.05$, \code{train\_bd\_disc.py:130--134}):
\begin{equation}
\boxed{\;\mathcal L_{\text{disc}}=\E\big[-\tilde y\log\sigma(\ell_\phi)-(1-\tilde y)\log\big(1-\sigma(\ell_\phi)\big)\big],\quad
\tilde y=\begin{cases}1-\epsilon&(\text{pos})\\ \epsilon&(\text{neg})\end{cases}\;}
\end{equation}
$\ell_\phi=\lambda\tanh(z_\phi/\lambda)$ (bounded logit, $\lambda{=}4$). 입력은
clean 부분경로 $[\xpre\Vert x_0]$ + 시나리오 인접성(시간 조건 없음).
positive = except 경로(무작위 절단 \code{make\_partial}); negative 3택:
\begin{center}\small
\begin{tabular}{@{}l l l@{}}
\toprule
\code{-neg} & negative 분포 & 최적해 $\sigma(\ell_\phi)$ 가 주는 비율\\
\midrule
\code{data} & normal 실제 경로 & $q_{\text{exc}}/(q_{\text{exc}}{+}q_{\text{norm}})\Rightarrow D/(1{-}D){=}q_{\text{exc}}/q_{\text{norm}}$\\
\code{model}& 그 블록 모델 unconditional 생성 & $\Rightarrow D/(1{-}D){=}q_{\text{exc}}/\pth$ \ (guidance 가 요구하는 분모)\\
\code{mix} & 50\% normal + 50\% model & 위 둘의 혼합\\
\bottomrule
\end{tabular}
\end{center}
\begin{remark}
Bayes-최적 판별자에서 $\sigma(\ell_\phi)=\frac{q}{q+p_{\text{neg}}}$ 이므로 $\frac{D}{1-D}=\frac{q}{p_{\text{neg}}}$.
\code{-neg model} 이면 $p_{\text{neg}}=\pth$ 라 guidance 가중치 $w=q/\pth$ 를 정확히 학습한다
(자매 문서 \code{BD\_BLOCK\_GUIDANCE\_MATH.pdf} 식~(3)·Lemma 1).
\end{remark}

\subsection*{Truncation(부분경로 길이) 표집과 불편성}
위 기댓값 $\E$ 는 \emph{두 단계} 표집을 포함한다: 전체 경로 $x$(길이 $L(x)$)를 $q$(또는
$p_{-}$)에서 뽑고, truncation kernel $T(j\mid x)$ 로 부분경로
$\tau_j(x)=[\text{dst}]\Vert(x_0,\dots,x_{j-1})$ 를 만든다(\code{make\_partial}). 명시하면
($\mathrm{BCE}_{\tilde y}(\ell)=-\tilde y\log\sigma(\ell)-(1-\tilde y)\log(1-\sigma(\ell))$,
$\tilde y_+{=}1-\epsilon,\ \tilde y_-{=}\epsilon$):
\begin{equation}
\mathcal L_{\text{disc}}
=\E_{x\sim q}\,\E_{j\sim T(\cdot\mid x)}\big[\mathrm{BCE}_{\tilde y_+}(\ell_\phi(\tau_j x))\big]
+\E_{x\sim p_{-}}\,\E_{j\sim T(\cdot\mid x)}\big[\mathrm{BCE}_{\tilde y_-}(\ell_\phi(\tau_j x))\big].
\end{equation}
$T$ 가 유도하는 부분경로 측도 $\mu_q(\pi)=\sum_x q(x)\sum_j T(j\mid x)\,\mathbf 1[\tau_j x{=}\pi]$
(neg 도 동일) 하에서 최적해는 $\sigma(\ell^\star_\phi(\pi))=\mu_q(\pi)/(\mu_q(\pi){+}\mu_{-}(\pi))$,
즉 $w=\tfrac{D}{1-D}=\mu_q(\pi)/\mu_{-}(\pi)$.

\textbf{불편성 조건.} guidance 가 요구하는 것은 \emph{clean} prefix 밀도비 $Q(\pi)/P_{-}(\pi)$
(여기서 $Q(\pi)=\sum_{x\to\pi}q(x)$: $\pi$ 로 시작하고 같은 dst 로 끝나는 경로들의 주변확률,
$x\to\pi$ 는 ``$x$ 가 어떤 절단으로 $\pi$ 를 생성''). 그런데
\begin{equation}
\frac{\mu_q(\pi)}{\mu_{-}(\pi)}
=\frac{\sum_{x\to\pi} q(x)\,T(j_\pi\mid x)}{\sum_{x\to\pi} p_{-}(x)\,T(j_\pi\mid x)}
\;=\;\frac{Q(\pi)}{P_{-}(\pi)}\quad\Longleftarrow\quad T(j\mid x)=\tau(\pi)\ \text{(완성경로 $x$ 에 무관)},
\end{equation}
즉 truncation 가중이 $\pi$ 에만 의존해야 분자·분모에서 상쇄되어 \emph{불편(unbiased)}이 된다.

\textbf{현재 구현 — per-path uniform (편향).} $T(j\mid x)=\dfrac{1}{L(x)-m+1}$
($m{=}$\code{min\_len}, \code{bd\_disc.py:167})는 완성경로 길이 $L(x)$ 에 의존한다
$\Rightarrow$ 상쇄되지 않아 $w$ 가 \textbf{완성경로 길이로 재가중된 편향 추정}이 된다
(긴 완성경로일수록 $1/(L{-}m{+}1)$ 로 낮게 가중). 게다가 부분경로 길이의 \emph{주변분포}가
경로길이 분포로 쏠려 짧은 prefix 과대·긴 prefix 과소표집.

\textbf{불편형 — global-uniform length.} 경로와 무관하게 $\ell\sim U\{m,\dots,\ell_{\max}\}$ 를
먼저 뽑고 $L(x)\ge\ell$ 인 경로만 $\tau_\ell(x)$ 로 자르면, nonzero 가중이
$\tfrac{1}{\ell_{\max}-m+1}$ 로 \textbf{상수($x$ 무관)} $\Rightarrow$
$\mu_q(\pi)=\tfrac{1}{\ell_{\max}-m+1}Q(\pi)$ 이므로 $w=Q(\pi)/P_{-}(\pi)$ (불편) + 길이
주변분포 평탄화.
\begin{remark}
단 $\ell$ 이 큰 구간은 ``그 길이 이상인'' 실제 경로 수에 상한이 걸려(대부분 15--25, 최대
$\approx 61$) 소수의 긴 경로를 재사용하게 된다 — tail 커버리지는 개선되나 데이터 희소성이
상한을 정한다. 따라서 global-uniform 은 (i) 길이 의존 편향 제거, (ii) tail 표집 완화의 두
효과를 주되, 긴 부분경로의 \emph{다양성}은 근본적으로 데이터에 묶인다.
\end{remark}

% =====================================================================
\section{판별기 학습 loss (2) — D-CBG classifier (노이즈 상태)}
% =====================================================================
D-CBG classifier 는 \emph{노이즈} 상태 $\xt$ 를 시간 $t$ 로 조건화하여 같은 BCE 로 학습한다
(\code{train\_dcbg\_classifier.py:114--119}):
\begin{equation}
\boxed{\;\mathcal L_{\text{clf}}=\E\big[-\tilde y\log\sigma\big(z_\phi(\xt,t)\big)-(1-\tilde y)\log\big(1-\sigma(z_\phi(\xt,t))\big)\big]\;}
\end{equation}
positive/negative 를 \textbf{생성 시점과 동일한 forward} 로 손상시켜 입력한다:
\begin{itemize}[nosep,leftmargin=1.4em]
  \item mask: \code{corrupt\_mask} — $j$ 개 확정 블록 + 현재 블록 rate $t$ 마스킹(\code{dcbg\_plugin.py:107}).
  \item graph: \code{corrupt\_graph} — 현재 블록을 $q(\xt\mid\xz)$ 로 정수 $t$ 노이징(\code{:132}).
\end{itemize}
따라서 $z_\phi$ 는 \textbf{전 노이즈레벨에 걸친} $p_\phi(y\mid \xt,t)$ 를 학습한다 — IW disc(§4)와의
유일한 차이는 이 \emph{노이즈-상태 + 시간 입력}이며, 아키텍처($\code{DCBGClassifierAdj}$: 동일 GraphSAGE
+ edge 피처 + bounded logit)·데이터·negative pool 은 통제된다.

% =====================================================================
\section{수식 $\leftrightarrow$ 코드 대응표}
% =====================================================================
\begin{center}
\renewcommand{\arraystretch}{1.2}
\begin{tabular}{@{}p{0.5\linewidth} l@{}}
\toprule
\textbf{수식} & \textbf{코드 위치}\\
\midrule
canvas·loss mask (길이 처리) & \code{bd\_models.py: build\_canvas, :435, :498}\\
mask SUBS $1/t$ 가중 CE $\mathcal L_{\text{mask}}$ & \code{bd\_models.py:437--464}\\
graph $q(\xt\mid\xz)$ 노이징 & \code{bd\_models.py:504--506}\\
true/pred posterior (D3PM) & \code{bd\_models.py:517--531}\\
$\mathcal L_{\text{KL}},\mathcal L_{\text{CE}},\mathcal L_{\text{con}}$ & \code{bd\_models.py:537--552}\\
ce-gating 총합 $\mathcal L_{\text{graph}}$ & \code{bd\_trainer.py:39--52}\\
IW disc BCE + label smoothing & \code{train\_bd\_disc.py:130--134}\\
negative 3택 (data/mix/model) & \code{train\_bd\_disc.py:94--128}\\
D-CBG classifier BCE (노이즈 입력) & \code{train\_dcbg\_classifier.py:114--119}\\
생성-정합 corruption & \code{dcbg\_plugin.py:107 / :132}\\
\bottomrule
\end{tabular}
\end{center}

\end{document}
